\documentclass[semifinal,survey]{cpecmu}
%% This is a sample document demonstrating how to use the CPECMU
%% project template. If you are having trouble, see "cpecmu.pdf" for
%% documentation.

\projectNo{S807-2}
\acadyear{2021}

\titleTH{โปรแกรมวางแผนปฏิทินการศึกษา มหาวิทยาลัยเชียงใหม่}
\titleEN{Chiang Mai University's academic calendar planner}

\author{นายเจษฎา จินะกะ}{Jetsada Jinaka}{620612144}
\author{นายเอื้อบุญ เรือนคำฟู}{Aueboon Ruanekamfu}{620612170}

\cpeadvisor{chinawat}
\cpecommittee{pruet}
\cpecommittee{lachana}

%% Some possible packages to include:
\usepackage[final]{graphicx} % for including graphics

%% Add bookmarks and hyperlinks in the document.
\PassOptionsToPackage{hyphens}{url}
\usepackage[colorlinks=true,allcolors=Blue4,citecolor=red,linktoc=all]{hyperref}
\def\UrlFont{\thaifonttt}

%% Set up commenting
\iffinal
  \usepackage[disabled]{authcomments}
\else
  \usepackage{authcomments}
\fi
\newcommenter{CI}{0.0,0.5625,0.0}  % green

%% Needed just by this example, but maybe not by most reports
\usepackage{afterpage} % for outputting
\usepackage{pdflscape} % for landscape figures and tables. 

%% Some other useful packages. Look these up to find out how to use
%% them.
% \usepackage{natbib}    % for author-year citation styles
% \usepackage{txfonts}
% \usepackage{appendix}  % for appendices on a per-chapter basis
% \usepackage{xtab}      % for tables that go over multiple pages
% \usepackage{subfigure} % for subfigures within a figure
% \usepackage{pstricks,pdftricks} % for access to special PostScript and PDF commands
% \usepackage{nomencl}   % if you have a list of abbreviations

%% if you're having problems with overfull boxes, you may need to increase
%% the tolerance to 9999
% \tolerance=9999

\bibliographystyle{plain}
% \bibliographystyle{IEEEbib}

% \renewcommand{\topfraction}{0.85}
% \renewcommand{\textfraction}{0.1}
% \renewcommand{\floatpagefraction}{0.75}

%% Example for glossary entry
%% Need to use glossary option
%% See glossaries package for complete documentation.
\ifglossary
  \newglossaryentry{lorem ipsum}{
    name=lorem ipsum,
    description={derived from Latin dolorem ipsum, translated as ``pain itself''}
  }
\fi

%% Uncomment this command to preview only specified LaTeX file(s)
%% imported with \include command below.
%% Any other file imported via \include but not specified here will not
%% be previewed.
%% Useful if your report is large, as you might not want to build
%% the entire file when editing a certain part of your report.
% \includeonly{chapters/intro,chapters/background}

\begin{document}
\maketitle
\makesignature

\ifproject
   \begin{abstractTH}
      \par ในแต่ละปีการศึกษา สำนักทะเบียนและประมวลผลมหาวิทยาลัยเชียงใหม่ จำเป็นจะต้องจะต้องจัดทำร่างปฏิทินการศึกษาสำหรับปีการศึกษาถัดไปเป็นประจำทุกปี \enskip
      ซึ่งในการทำร่างปฏิทิน ซึ่งในปฏิทินประกอบไปด้วยกำหนดการของกิจกรรมการศึกษาต่างๆ เช่น วันเปิดภาคเรียน วันลงทะเบียนเรียน วันสอบ เป็นต้น\enskip
      ซึ่งวันกำหนดการต่างๆ จะมีเงื่อนไขที่อ้างอิงกับวันกำหนดการต่างๆ ซึ่งทำให้การทำร่งปฏิทินขึ้นมาหนึ่งร่าง มีความยุ่งยากในการทำมากจนใช้เวลาเป็นสัปดาห์\enskip
      และยังต้องมาแก้ใหม่หากกรรมการของมหาวิทยาลัยเชียงใหม่ไม่ให้ความเห็นชอบ\enskip

      จากปัญหาข้างต้น ผู้จัดทำได้จัดทำโปรแกรมวางแผนปฏิทินการศึกษา มหาวิทยาลัยเชียงใหม่ เพื่อเป็นเครื่องมือที่ช่วยให้ผู้จัดทำร่างปฏิทินของสำนักทะเบียน
      เพื่อช่วยลดระยะในการจัดทำปฏิทินลงได้ อีกทั้งสามารถเพิ่มวามสะดวกในการจัดทำร่างปฏิทินได้ สามารถแก้ไข สร้าง หรือทำซ้ำร่างปฏิทินที่มีฉบับที่คล้าคลึงกัน
      โปรแกรมดังกล่าวก็สามารถตอบสนองได้
   \end{abstractTH}

   \begin{abstract}
      in each academic year Chiang Mai University Registration and Processing Office It is necessary to make
      Outline an academic calendar for the next academic year annually. which in drafting the calendar in which the accompanying calendar
      along with the schedule of various educational activities such as the opening day of the semester registration day, exam date, etc.
      which scheduled dates There will be conditions that refer to various scheduled dates. which makes one draft of the calendar have
      The complexity of making it took weeks. And still have to come and fix it again if the director of the University of Chiang-
      new disapproval
      from the above problem The organizer has prepared an academic calendar planning program. Chiang Mai University for
      A tool that helps drafters of the registry's calendar. To help reduce the time for making calendars, it can also be
      It is convenient to create a draft calendar. You can edit, create or duplicate a draft calendar with similar editions.
      together, such programs can respond.
   \end{abstract}

   \iffalse
      \begin{dedication}
         This document is dedicated to all Chiang Mai University students.

         Dedication page is optional.
      \end{dedication}
   \fi % \iffalse

   \begin{acknowledgments}
           โครงการนี้รับความกรุณาจาก อ.ดร.ชินวัตรอิศราดิสัยกุล อาจารย์ที่ปรึกษาโครงการ
      ที่คอยให้ความช่วยเหลือ เสนอแนะแนวทางในการแก้ปัญหา คอยติดตามงาน และจัดลําดับความสําคัญงานอยู่เสมอ
      รวมไปถึงขอบคุณอาจารย์คณะกรรมการทั้ง อ.ดร. พฤษภ์ บุญมา และ ผศ.ดร. ลัชนา ระมิงค์วงศ์ ที่ให้คำแนะนำต่างๆ รวมไปถึงเพื่อนๆที่ให้คอยให้กำลังใจและช่วยเหลือ ให้คำแนะนำมาตลอด
      ทำให้โครงงานเล่มนี้ออกมาได้อย่างสมบูรณ์ 
      ขอขอบคุณภาควิชาวิศวกรรมคอมพิวเตอร์ คณะวิศวกรรมศาสตร์ มหาวิทยาลัยเชียงใหม่ ที่มอบทุนสนับสนุนการพัฒนา และเอื้อเฟื้อสถานที่ในการจัดทําโครงการนี้
      \acksign{2023}{3}{30}
   \end{acknowledgments}%
\fi % \ifproject

\contentspage

\ifproject
   \figurelistpage % สารบัญรูป 

   \tablelistpage  % สารบัญตาราง
\fi % \ifproject

% \abbrlist % this page is optional

% \symlist % this page is optional

% \preface % this section is optional


\pagestyle{empty}\cleardoublepage
\normalspacing \setcounter{page}{1} \pagenumbering{arabic} \pagestyle{cpecmu}

\chapter{\ifenglish Introduction\else บทนำ\fi}

\section{\ifenglish Project rationale\else ที่มาของโครงงาน\fi}
    ในปัจจุบันปฏิทินการศึกษาของมหาวิทยาลัยเชียงใหม่นั้นจัดทำโดยทางสำนักทะเบียนของมหาวิทยาลัยเชียงใหม่ 
โดยที่ในการที่จะสร้างปฏิทินขึ้นมาได้นั้น จะต้องมีการร่างโครงร่างของปฏิทินออกมา โดยที่ในการที่จะร่างปฏิทินนั้ก็จะต้องมีเงื่อนไขในการส้รางปฏิทินต่างๆ ไม่ว่าจะเป็น การกำหนดวันเปิดภาคเรียน การกำหนดระยะเวลาของคาบเรียน วันปิดภาคเรียน ไปจนถึงวันลงทะเบียน 
เราได้พบปัญหาว่าใน 1 ปฏิทินการศึกษานั้น ใช้เวลาในการร่างปฏิทินไม่ต่ำกว่า 2 สัปดาห์ เนื่องจากทางผู้จัดทำต้องกำหนดเงื่อนไขที่ได้กล่าวข้างต้น ถึงจะทำการวางโครงร่างของปฏิทินได้
และเมื่อสามารถวางโครงร่างได้แล้วก็จะต้องนำโครงร่างไปส่งให้เป็นที่พิจารณาแก่คณะกรรมการบริหารของมหาวิทยาลัยเชียงใหม่ ซึ่งจะทำให้ต้องมีการแก้ไขอยู่หลายๆครั้ง จึงทำให้การทำปฏิทินวันเปิดภาคเรียนนั้นใช้เวลานานมากเกินไป เราจึงได้เล็งปัญหาของการสร้างปฏทินการศึกษา
ซึ่งจัดทำโดยสำนักทะเบียนมหาวิทยาลัยเชียงใหม่ จึงได้เกิดเป็น


\section{\ifenglish Objectives\else วัตถุประสงค์ของโครงงาน\fi}
\begin{enumerate}
    \item 
    \item
    \item
    \item
    \item

\end{enumerate}

\section{\ifenglish Project scope\else ขอบเขตของโครงงาน\fi}

\subsection{\ifenglish Hardware scope\else ขอบเขตด้านฮาร์ดแวร์\fi}

\subsection{\ifenglish Software scope\else ขอบเขตด้านซอฟต์แวร์\fi}

\section{\ifenglish Expected outcomes\else ประโยชน์ที่ได้รับ\fi}

\section{\ifenglish Technology and tools\else เทคโนโลยีและเครื่องมือที่ใช้\fi}

\subsection{\ifenglish Hardware technology\else เทคโนโลยีด้านฮาร์ดแวร์\fi}

\subsection{\ifenglish Software technology\else เทคโนโลยีด้านซอฟต์แวร์\fi}

\section{\ifenglish Project plan\else แผนการดำเนินงาน\fi}

\begin{plan}{6}{2020}{2}{2021}
    \planitem{7}{2020}{8}{2020}{ศึกษาค้นคว้า}
    \planitem{8}{2020}{1}{2021}{ชิล}
    \planitem{2}{2021}{2}{2021}{เผา}
    \planitem{12}{2019}{1}{2022}{ทดสอบ}
\end{plan}

\section{\ifenglish Roles and responsibilities\else บทบาทและความรับผิดชอบ\fi}
อธิบายว่าในการทำงาน นศ. มีการกำหนดบทบาทและแบ่งหน้าที่งานอย่างไรในการทำงาน จำเป็นต้องใช้ความรู้ใดในการทำงานบ้าง

\section{\ifenglish%
Impacts of this project on society, health, safety, legal, and cultural issues
\else%
ผลกระทบด้านสังคม สุขภาพ ความปลอดภัย กฎหมาย และวัฒนธรรม
\fi}

แนวทางและโยชน์ในการประยุกต์ใช้งานโครงงานกับงานในด้านอื่นๆ รวมถึงผลกระทบในด้านสังคมและสิ่งแวดล้อมจากการใช้ความรู้ทางวิศวกรรมที่ได้

\chapter{\ifenglish Background Knowledge and Theory\else ทฤษฎีที่เกี่ยวข้อง\fi}

การทำโครงงาน เริ่มต้นด้วยการศึกษาค้นคว้า ทฤษฎีที่เกี่ยวข้อง หรือ งานวิจัย/โครงงาน ที่เคยมีผู้นำเสนอไว้แล้ว ซึ่งเนื้อหาในบทนี้ก็จะเกี่ยวกับการอธิบายถึงสิ่งที่เกี่ยวข้องกับโครงงาน เพื่อให้ผู้อ่านเข้าใจเนื้อหาในบทถัดๆ ไปได้ง่ายขึ้น

\section{หน้าเว็บ}
หน้าเว็บ คือ หน้าเอกสารที่ถูกแสดงโดย เว็บเบราว์เซอร์ เพ่ือแสดงข้อมูลต่างๆ ที่เป็นข้อความ รูปภาพ และส่ือผสมต่างๆ ซึ่งเนื้อหาของหน้าเว็บเป็นอย่างไร ขึ้นอยู่กับวัตถุประสงค์ของ เจ้าของหน้าเว็บ ไม่ว่าจะเป็นเนื้อหาเกี่ยวกับการศึกษาธุรกิจ หรือ ความบันเทิง เป็นต้น


\section{ HTML}
HTML ย่อมาจาก HyperText Markup Language เป็น ภาษาคอมพิวเตอร์ท่ีใช้สร้างหน้าเว็บในรูปแบบของ ไฟล์ HTML (คือไฟล์ที่มีนามสกุลเป็น .htm หรือ .html) ซึ่งมีเว็บเบราว์เซอร์เป็นโปรแกรมที่ใช้แปลงไฟล์ HTML เพื่อ แสดงผลในรูปของหน้าเว็บ
ไฟล์ HTML เป็นไฟล์รหัสแอสกี (ASCII) ถูกบันทึกในรูปของไฟล์เอกสาร (Text File) ที่สามารถถูกสร้างจากโปรแกรมสร้างไฟล์ ข้อความ เช่น Notepad หรือ Word Processing ทั่วๆ ไป ซึ่งลักษณะของไฟล์ HTML ประกอบไปด้วยแท็กต่างๆ ที่เป็นคําาส่ังของ HTML ซึ่งแท็กจะอยู่ภายในเครื่องหมาย <และ>

\section{ CSS}
CSS ย่อมาจาก Cascading Style Sheet  มักเรียกโดยย่อว่า "สไตล์ชีต" คือภาษาที่ใช้เป็นส่วนของการจัดรูปแบบการแสดงผลเอกสาร  HTML โดยที่ CSS กำหนดกฏเกณฑ์ในการระบุรูปแบบ (หรือ "Style") ของเนื้อหาในเอกสาร 
อันได้แก่ สีของข้อความ สีพื้นหลัง ประเภทตัวอักษร และการจัดวางข้อความ ซึ่งการกำหนดรูปแบบ หรือ Style นี้ใช้หลักการของการแยกเนื้อหาเอกสาร HTML ออกจากคำสั่งที่ใช้ในการจัดรูปแบบการแสดงผล กำหนดให้รูปแบบของการแสดงผลเอกสาร 
ไม่ขึ้นอยู่กับเนื้อหาของเอกสาร เพื่อให้ง่ายต่อการจัดรูปแบบการแสดงผลลัพธ์ของเอกสาร HTML

\section{JavaScript}
JavaScript หรือย่อด้วย JS เป็นภาษาเขียนโปรแกรมที่ถูกพัฒนาและปฏิบัติตามข้อกำหนดมาตรฐานของ ECMAScript
ภาษา JavaScript นั้นเป็นภาษาระดับสูง คอมไพล์ในขณะที่โปรแกรมรัน (JIT) และเป็นภาษาเขียนโปรแกรมแบบหลายกระบวนทัศน์ เช่น การเขียนโปรแกรมเชิงขั้นตอน การเขียนโปรแกรมเชิงวัตถุ หรือการเขียนโปรแกรมแบบ Functional; 
ภาษา JavaScript มีไวยากรณ์ที่เหมือนกับภาษา C ใช้วงเล็บเพื่อกำหนดบล็อคของคำสั่ง นอกจากนี้ JavaScript ยังเป็นภาษาที่มีประเภทข้อมูลแบบไดนามิกส์ เป็นภาษาแบบ Prototype-based และ First-class function

\section{Node.JS}
Node.js คือสภาพแวดล้อมการทำงานของภาษา JavaScript นอกเว็บเบราว์เซอร์ที่ทำงานด้วย V8 engine นั่นคือสามารถใช้ Node.js ในการพัตนาแอพพลิเคชันแบบ Command line แอพพลิเคชัน Desktop หรือแม้แต่เว็บเซิร์ฟเวอร์ได้ 
โดยที่ Node.js จะมี APIs ที่เราสามารถใช้สำหรับทำงานกับระบบปฏิบัติการ เช่น การรับค่าและการแสดงผล การอ่านเขียนไฟล์ และการทำงานกับเน็ตเวิร์ก เป็นต้น

\section{React.JS}
React เป็น JavaScript library ที่ใช้สำหรับสร้าง user interface ในฝั่งด้าน Front-end ที่ให้เราสามารถเขียนโค้ดในการสร้าง UI ที่มีความซับซ้อนแบ่งเป็นส่วนเล็กๆออกจากกันได้ ซึ่งแต่ละส่วนสามารถแยกการทำงานออกจากกันได้อย่างอิสระ 
และทำให้สามารถนำชิ้นส่วน UI เหล่านั้นไปใช้ซ้ำได้อีก

\section{MongoDB}
MongoDB เป็น open-source document database โดยเป็นฐานข้อมูลแบบ NoSQL หรือเรียกว่าไม่มีความสัมพันธ์ของตารางแบบ SQL ทั่วๆไป แต่จะเก็บข้อมูลเป็นแบบ JSON แทน การบันทึกข้อมูลทุกๆ record ใน MongoDB 
เราจะเรียกมันว่า Documentซึ่งจะเก็บค่าเป็น key และ value หรือก็คือไฟล์ JSON

\section{JSON}
ย่อมาจาก JavaScript Object Notation เป็นฟอร์แมตสำหรับแลกเปลี่ยนข้อมูลคอมพิวเตอร์ ฟอร์แมต JSON นั้นอยู่ในรูปข้อความธรรมดา (plain text) ที่ทั้งมนุษย์และโปรแกรมคอมพิวเตอร์สามารถอ่านเข้าใจได้ 
มาตรฐานของฟอร์แมต JSON คือ RFC 4627 มี Internet media type เป็น application/json และมีนามสกุลของไฟล์เป็น .json
ปัจจุบัน JSON นิยมใช้ในเว็บแอปพลิเคชัน โดยเฉพาะ AJAX โดย JSON เป็นฟอร์แมตทางเลือกในการส่งข้อมูล นอกเหนือไปจาก XML ซึ่งนิยมใช้กันอยู่แต่เดิม สาเหตุที่ JSON เริ่มได้รับความนิยมเป็นเพราะกระชับและเข้าใจง่ายกว่า XML

JSON นั้นใช้ความสัมพันธ์ของภาษาจาวาสคริปต์ แต่ไม่ถูกมองว่าเป็นภาษาโปรแกรม กลับถูกมองว่าเป็นภาษาในการแลกเปลี่ยนข้อมูลมากกว่า ในปัจจุบันมีไลบรารีของภาษาโปรแกรมอื่นๆ ที่ใช้ประมวลผลข้อมูลในรูปแบบ JSON มากมาย







%\section{Third section}
%Section 3 text. The dielectric constant\index{dielectric constant}
%at the air-metal interface determines
%the resonance shift\index{resonance shift} as absorption or capture occurs
%is shown in Equation~\eqref{eq:dielectric}:

%\begin{equation}\label{eq:dielectric}
%k_1=\frac{\omega}{c({1/\varepsilon_m + 1/\varepsilon_i})^{1/2}}=k_2=\frac{\omega
%\sin(\theta)\varepsilon_\mathit{air}^{1/2}}{c}
%\end{equation}

%\noindent
%where $\omega$ is the frequency of the plasmon, $c$ is the speed of
%light, $\varepsilon_m$ is the dielectric constant of the metal,
%$\varepsilon_i$ is the dielectric constant of neighboring insulator,
%and $\varepsilon_\mathit{air}$ is the dielectric constant of air.

%\section{About using figures in your report}

% define a command that produces some filler text, the lorem ipsum.
%\newcommand{\loremipsum}{
%  \textit{Lorem ipsum dolor sit amet, consectetur adipisicing elit, sed do
%  eiusmod tempor incididunt ut labore et dolore magna aliqua. Ut enim ad
%  minim veniam, quis nostrud exercitation ullamco laboris nisi ut
%  aliquip ex ea commodo consequat. Duis aute irure dolor in
%  reprehenderit in voluptate velit esse cillum dolore eu fugiat nulla
%  pariatur. Excepteur sint occaecat cupidatat non proident, sunt in
%  culpa qui officia deserunt mollit anim id est laborum.}\par}

%\begin{figure}
%  \centering

%  \fbox{
%     \parbox{.6\textwidth}{\loremipsum}
%  }

  % To include an image in the figure, say myimage.pdf, you could use
  % the following code. Look up the documentation for the package
  % graphicx for more information.
  % \includegraphics[width=\textwidth]{myimage}

%  \caption[Sample figure]{This figure is a sample containing \gls{lorem ipsum},
%  showing you how you can include figures and glossary in your report.
%  You can specify a shorter caption that will appear in the List of Figures.}
%  \label{fig:sample-figure}
%\end{figure}

%Using \verb.\label. and \verb.\ref. commands allows us to refer to
%figures easily. If we can refer to Figures
%\ref{fig:walrus} and \ref{fig:sample-figure} by name in the {\LaTeX}
%source code, then we will not need to update the code that refers to it
%even if the placement or ordering of the figures changes.

%\loremipsum\loremipsum

% This code demonstrates how to get a landscape table or figure. It
% uses the package lscape to turn everything but the page number into
% landscape orientation. Everything should be included within an
% \afterpage{ .... } to avoid causing a page break too early.
%\afterpage{
%  \begin{landscape}
%  \begin{table}
%    \caption{Sample landscape table}
%    \label{tab:sample-table}

%   \centering

%    \begin{tabular}{c||c|c}
%        Year & A & B \\
%        \hline\hline
%        1989 & 12 & 23 \\
%        1990 & 4 & 9 \\
%        1991 & 3 & 6 \\
%    \end{tabular}
%  \end{table}
%  \end{landscape}
%}

%\loremipsum\loremipsum\loremipsum

%\section{Overfull hbox}

%When the \verb.semifinal. option is passed to the \verb.cpecmu. document class,
%any line that is longer than the line width, i.e., an overfull hbox, will be
%highlighted with a black solid rule:
%\begin{center}
%\begin{minipage}{2em}
%juxtaposition
%\end{minipage}
%\end{center}

%\section{\ifenglish%
%\ifcpe CPE \else ISNE \fi knowledge used, applied, or integrated in this project
%\else%
%ความรู้ตามหลักสูตรซึ่งถูกนำมาใช้หรือบูรณาการในโครงงาน
%\fi
%}

%อธิบายถึงความรู้ และแนวทางการนำความรู้ต่างๆ ที่ได้เรียนตามหลักสูตร ซึ่งถูกนำมาใช้ในโครงงาน

%\section{\ifenglish%
%Extracurricular knowledge used, applied, or integrated in this project
%\else%
%ความรู้นอกหลักสูตรซึ่งถูกนำมาใช้หรือบูรณาการในโครงงาน
%\fi
%}

%อธิบายถึงความรู้ต่างๆ ที่เรียนรู้ด้วยตนเอง และแนวทางการนำความรู้เหล่านั้นมาใช้ในโครงงาน

\chapter{\ifproject%
\ifenglish Project Structure and Methodology\else โครงสร้างและขั้นตอนการทำงาน\fi
\else%
\ifenglish Project Structure\else โครงสร้างของโครงงาน\fi
\fi
}

ในบทนี้จะกล่าวถึงหลักการ และการออกแบบระบบไปจนถึงขั้นตอนการออกแบบจากความต้องการของผู้ใช้งาน

\makeatletter

% \renewcommand\section{\@startsection {section}{1}{\z@}%
%                                    {13.5ex \@plus -1ex \@minus -.2ex}%
%                                    {2.3ex \@plus.2ex}%
%                                    {\normalfont\large\bfseries}}

\makeatother
%\vspace{2ex}
% \titleformat{\section}{\normalfont\bfseries}{\thesection}{1em}{}
% \titlespacing*{\section}{0pt}{10ex}{0pt}

\section{การติดต่อและคุยงานเพื่อสรุปความต้องการของสำนักทะเบียน}

% \begin{figure}
% \begin{center}
% % \includegraphics{800px-Briny_Beach.jpg}
% \end{center}
% % \caption[Poem]{The Walrus and the Carpenter}
% \label{fig:walrus}
% \end{figure}

\section{เงื่อนไขของการวางร่างปฏิทินการศึกษา}
  เนื่องจากจุดประสงค์ของโครงงานนี้คือต้องการพัฒนาเว็บไซต์ให้แก่สำนักทะเบียน
  จึงจะต้องเริ่มจากการพูดคุยกับบุคลากรของสำนักทะเบียนเพื่อให้ได้ความต้องการที่แท้จริงของโครงงาน โดยในปฏิทินจะมีเงื่อนไขต่างๆ อันสรุปได้ดังนี้

\subsection{ภาคการศึกษาที่ 1}
  คาบเรียนของแต่ละวัน มีจำนวนดังนี้ วันจันทร์อย่างน้อย 14 คาบ วันอังคารอย่างน้อย 14 คาบ วันพุธอ่างน้อย 14 คาบ วันพฤหัสบดี 13 คาบ วันศุกร์ 12 คาบ(ประมาณ)\CIreply{is this from just one academic year, or from a requirement?}
\begin{itemize}
  \item วันเปิดภาคเรียนมักจะอยู่ในเดือนมิถุนายน
  \item รูปแบบวันจันทร์ พฤหัสบดี มีระยะเวลาการเรียนการสอนตลอดภาคการศึกษา จำนวน 42 ชั่วโมง
\end{itemize}
  \CIreply{convert to itemize, as shown in the first few items above}
  - รูปแบบวันอังคาร ศุกร์ มีระยะเวลาการเรียนการสอนตลอดภาคการศึกษา จำนวน 40 ชั่วโมง 30 นาที\CIreply{why?} \\
  - วันลงทะเบียนเรียนล่วงหน้า สัปดาห์แรกของเดือนก่อนหน้าที่จะเปิดภาคการศึกษา \\
  - ประกาศผลการลงทะเบียนเรียนล่วงหน้าหลังจากลงทะเบียนล่วงหน้าประมาณ 10 วัน \\
  - วันลงทะเบียนเรียนมักจะเริ่มก่อนวันเปิดภาคเรียนวันแรก 1--2 วัน \\
  - วันลงทะเบียนมีระยะเวลา 8 วัน (เป็นเวลาที่ตรงกับวันเรียนด้วย) \\
  - วันถอนกระบวนวิชาโดยไม่ได้รับอักษรลำดับขั้น W จะเริ่มจากวันลงทะเบียนเรียน \CI{ถึงวันที่ 1 ของเดือนถัดไป}{really?} \\
  - วันที่อาจารย์ที่ปรึกษาให้ความเห็นชอบการลงทะเบียนเรียนของนักศึกษาทางระบบ Internet ระยะเวลาเริ่มจาก 1 สัปดาห์หลังจากเปิดภาคเรียน\CI{จนถึงวันที่ 1 ของเดือนถัดไป}{really?} \\
  - หลังจากหมดวันถอนกระบวนวิชาโดยไม่ได้รับอักษรลำดับขั้น W จะเริ่มนับวันถอนกระบวนวิชาโดยได้รับอักษรลำดับขั้น W \\
  - หลังจากเปิดภาคเรียน เมื่อเรียนครบ 8 สัปดาห์จะเริ่มสอบกลางภาค \\
  - สอบกลางภาคสัปดาห์ที่ 9 ของภาคเรียน มีระยะเวลาทั้งหมด 7 วัน\CIreply{what's different from last item?} \\
  - หลังจากวันสุดท้ายของการสอบกลางภาค เริ่มเรียนครึ่งภาคเรียนหลัง (สัปดาห์ที่ 10) \\
  - ครึ่งภาคเรียนหลังเวลาเรียนจะมีระยะเวลา 7 สัปดาห์ \CI{เริ่ม}{?}สอบกลางภาค \\
  - วันสอบปลายภาคมีระยะเวลา 2 สัปดาห์ \\
  - หลังสอบปลายภาคจะทำการหยุดเรียนจนจบสัปดาห์นั้นและหยุดเรียนเพิ่มอีก 1 สัปดาห์ และเปิดภาคเรียนที่2วนวันจันทร์ของสัปดาห์ถัดมา

  \subsection{ภาคการศึกษาที่ 2}
  คาบเรียนของแต่ละวัน มีจำนวนดังนี้ วันจันทร์ 12 คาบ วันอังคาร 15 คาบ วันพุธ 15 คาบ วันพฤหัสบดี 15 คาบ วันศุกร์ 14 คาบ \\
  - รูปแบบวันจันทร์ พฤหัสบดี มีระยะเวลาการเรียนการสอนตลอดภาคการศึกษา จำนวน 40 ชั่วโมง 30 นาที \\
  - รูปแบบวันอังคาร ศุกร์ มีระยะเวลาการเรียนการสอนตลอดภาคการศึกษา จำนวน 43 ชั่วโมง 30 นาที \\
  - วันลงทะเบียนเรียนล่วงหน้า \CI{ก่อนเปิดภาคเรียน 25 วัน}{really?} \\
  - ประกาศผลการลงทะเบียนล่วงหน้า \CI{1 สัปดาห์}{from when?} \\
  - เปิดภาคการศึกษาที่ 2 (ระยะเวลา 18 สัปดาห์เหมือนภาคการศึกษาที่ 1) \\
  - วันลงทะเบียนเรียนมักจะเริ่มก่อนวันเปิดภาคเรียนวันแรก 1--2 วัน \\
  - วันลงทะเบียนมีระยะเวลา8วัน (เป็นเวลาที่ตรงกับวันเรียนด้วย) \\
  - หลังจากสอบปลายภาคภาคเรียนที่2 จะหยุดเรียน3 สัปดาห์ \\
  - หลังจากหยุดเรียน3สัปดาห์ จะเริ่มเปิดภาคเรียนฤดูร้อน

  \subsection{ภาคการศึกษาฤดูร้อน}
  ภาคฤดูร้อน มีรูปแบบการเรียนการสอนตั้งแต่วันจันทร์--ศุกร์ ระยะเวลาที่ใช้สำหรับการเรียนการสอนแต่ละกระบวนวิชามีเวลาเรียนจำนวน 29 คาบ เท่ากับ 43 ชั่วโมง 30 นาที \\
  - ลงทะเบียนล่วงหน้าของภาคฤดูร้อนเริ่มหลังจากปิดภาคเรียนที่ 2 ไปแล้ว 1 สัปดาห์(ระยะเวลา 4 วัน) \\
  - ประกาศผลการลงทะเบียนล่วงหน้าหลังจากปิดการลงทะเบียน 2 วัน \\
  - วันเพิ่ม-ถอนกระบวนวิชา/ลงทะเบียน สำหรับนักศึกษาทุกระดับ ก่อนวันเปิดภาคเรียน 1 สัปดาห์ \\
  - ภาคฤดูร้อนระยะเวลาเรียน 6 สัปดาห์ \\
  - สัปดาห์ที่ 7 ของภาคฤดูร้อนจะเป็นสอบปลายภาคฤดูร้อน \\
  - หลังจากสอบปลายภาคฤดูร้อนแล้วจะปิดปีการศึกษา \\
  \CI{-วันไห้ครู หยุดเรียน}{why is this part of summer semester?} \\

\CI{หลังจากที่ได้เงื่อนไขทั้งหมดครบแล้ว จะนำเงื่อนไขเหล่านี้มาแยกออกจากระบบที่ตอบสนอง
และเพิ่มความสะดวกสบายของผู้ใช้}{what's the point of this paragraph?}

\section{การทำงานของโปรแกรม}
จากการสรุปความต้องการของสำนักทะเบียนผ่านทางบุคลากร จึงสรุปออกมาเป็น user flow หรือสิ่งที่แสดงเส้นทางของผู้ใช้แอพพลิเคชั่นได้ดังรูปที่~\ref{fig:user-flow-new}
\CIreply{this figure doesn't outline the entire user flow; did you mean to say something different?}
%
\begin{figure}[h]
\centering
\includegraphics[width=1\textwidth]{pic3.1.jpg}
\caption{การทำงานเมื่อผู้ใช้ต้องการสร้างแบบปฏิทินใหม่}
\label{fig:user-flow-new}
\end{figure}
%
โดยผู้ใช้จะเริ่มจากการเข้าสู่ระบบโดยใช้ CMU \CI{QAuth}{?} หลังจากนั้น คลิกที่สร้างดราฟใหม่ 
หลังจากนั้นเว็บไซต์จะต้องการทราบวันแรกของการเปิดภาคเรียนเพื่อนำไปสร้างปฏิทินการศึกษา
โดยหน้า Document จะเป็นหน้าที่ใช้จัดการกับร่างปฏิทินทั้งหมดที่ผู้ใช้ได้สร้างไว้
หลังจากที่ผู้ใช้ได้คลิกสร้างปฏิทินขึ้นมาใหม่ ระบบจะต้องการให้ผู้ใช้กรอกข้อมูลของวันเปิดเทอม
ของปีการศึกษานั้น หลังจากนั้นระบบจะทำการสร้างร่างปฏิทินการศึกษาแบบอัตโนมัติ
เพื่อทำให้ง่ายต่อการแก้ไข ไม่เกิดความยุ่งยากในการต้องมาเพิ่มกิจกรรมทีละวันกิจกรรม
โดยกิจกรรมที่นำไปใส่ลงในปฏิทินแบบอัติโนมัตินั้นจะได้มาจากการคำนวณ
วันที่อยู่ห่างจากวันเปิดเทอมตามเงื่อนไขของปฏิทิน
% วางรูป Dupilicate flow\
% รูปที่3.2 การทำงานเมื่อผู้ใช้ต้องการทำซ้ำปฏิทินเดิม

\section{โครงสร้างของโครงงาน}

\begin{figure}[h]
\centering
\includegraphics[width=1\textwidth]{pic3.2.jpg}
\caption{การทำงานเมื่อผู้ใช้ต้องการสร้างแบบปฏิทินใหม่}
\end{figure}

จากรูปที่~\CI{3.2}{use reference} ผู้ใช้ต้องการจะทำซ้ำ หลังจากที่ผู้ใช้อยู่ในหน้า Document และคลิกที่ทำซ้ำหรือ Duplicate
เว็บไซต์จะต้องการให้ผู้ใช้กรอกชื่อของปฏิทินที่จะสร้างใหม่ที่ทำซ้ำมาจากปฏิทินเดิม
และปีที่ต้องการเปลี่ยนใหม่ หากมีกิจกรรมของปฏิทินเดิมที่คล้ายคลึงกับปฏิทินของปีถัดไป 
ผู้ใช้สามารถทำซ้ำปฏิทินเดิมแล้วเปลี่ยนเป็นปีถัดไปได้เลย
\CIreply{what's different from the previous flow?}

\section{โครงสร้างโปรแกรม}
ในส่วนของ client จะใช้ภาษา React.JS ในการสร้างเว็บไซต์ แพลตฟอร์มนี้จะใช้กับ
คอมพิวเตอร์ โดยมี Node.JS ในส่วน backend และใช้ API ในการรับส่งกับฐานข้อมูล
และฐานข้อมูล MongoDB

\chapter{\ifproject%
\ifenglish Experimentation and Results\else การทดลองและผลลัพธ์\fi
\else%
\ifenglish System Evaluation\else การประเมินระบบ\fi
\fi}
\section{ทดสอบการลงชื่อเข้าใช้} 
   การทดสอบเพื่อ ตรวจสอบว่าผู้ที่เข้าใช้งานเว็บไซต์ เป็นบุคลากรของสำนักทะเบียน มหาวิทยาลัยเชียงใหม่ 
และเมื่อผู้เข้าใช้งานต้องการที่จะเข้าสู่ระบบ log in สามารถเข้าถึงเว็บไซต์ได้ตามปกติ

\section{ทดสอบความแม่นยำและประสิทธิภาพการใช้งานของการวางร่างปฏิทินแบบอัตโนมัติ}
   นำปฏิทินการศึกษาของปีที่ผ่านมามาทดสอบ โดยทดสอบให้ผู้ใช้สร้างปฏิทินการศึกษาโดยนำข้อมูลของปฏิทินจากปีที่ผ่านมา 
เพื่อให้ร่างปฏิทินการศึกษาโดยมีการกำหนดวันกิจกรรมแบบอัตโนมัติ แล้วจึงนำมาเปรียบเทียบกับปฏิทินการศึกษาที่เคยร่างไว้แล้ว 
เมื่อเปรียบเทียบกันแล้วผลลัพธ์ของการเปรียบเทียบที่ได้ปฏิทินการศึกษาที่จัดการโดยอัติโนมัติจะต้องไม่ด้อยไปกว่าปฏิทินการศึกษาที่เคยร่างไว้
โดยผู้ที่เป็นคนตัดสินใจว่าไม่ด้อยกว่านั้นเป็นผู้ที่จะมาใช้งานเว็บไซต์ หรือก็คือบุคลากรของสำนักทะเบียน\CIreply{เทียบกับอะไร?}

\section{การประเมินผลระบบ UX/UI } 
จะมีเกณฑ์หลักๆ ตามหลักของ Product ที่ดีดังนี้
\begin{enumerate}
   \item Usable การใช้ง่าย: โดยการนำเว็บไซต์นี้มาลองให้ผู้ใช้งานใช้และถาม feedback หลังใช้งาน
   \item Equitable การเท่าเทียมกัน: คือการนำเว็บไซต์นี้ไปลองใช้งานในคอมพิวเตอร์หลายๆ ขนาด\CIreply{จำเป็นไหม เพราะอะไร}
   \item Enjoyable: ความพึงพอใจของผู้ใช้มากกว่าการทำปฏิทินการศึกษาด้วยวิธีเดิมที่เคยใช้
   \item Useful มีประโยชน์: เว็บไซต์จะต้องแก้ปัญหาของผู้ใช้ตรงตามจุดประสงค์ของโครงงาน (เป็นการทดสอบการทำงานของระบบด้วย)
\end{enumerate}
\CIreply{performance evaluation?}

\ifproject
\include{chapters/conclusion}
\fi

\CIreply{use BibTeX entries properly!}
\bibliography{myReport}
% \begin{thebibliography}{1}
   
%   \bibitem{ Jirayut Intachai}
%   Jirayut Intachai. (2565). HTML คืออะไร? ทำไมคนเขียนเว็บไซต์ต้องใช้งาน. สืบค้น 16 กุมภาพันธ์ 2565,
%   จาก https://goterrestrial.com/2021/05/19/what-is-html/

%   \bibitem{ Kongvut Sangkla}
%   Kongvut Sangkla. (2565). ทำไมควรเรียนรู้ภาษา JavaScript ?. สืบค้น 16 กุมภาพันธ์ 2565,
%   จาก https://blog.2my.xyz/2020/03/10/javascript-from-zero-to-hero/

%   \bibitem{ Meta Platforms}
%   Meta Platforms. (2565).  React. สืบค้น 16 กุมภาพันธ์ 2565,
%   จาก https://reactjs.org/docs/getting-started.html
  
%   \bibitem{ Narator Game}
%   Narator Game. (2565). MongoDB คืออะไร?. สืบค้น 16 กุมภาพันธ์ 2565,
%   จาก https://medium.com/@narator.game/mongodb-คืออะไร-5a0844d7a3c8
  
%   \bibitem{ OpenJS Foundation}
%   OpenJS Foundation. (2565).  About Node.js®. สืบค้น 16 กุมภาพันธ์ 2565,
%   จาก https://nodejs.org/en/about/
  
%   \bibitem{Mitsumasa}
%   Mitsumasa. (2565). JSON คืออะไร. สืบค้น 16 กุมภาพันธ์ 2565,
%   จาก http://www.tutorialdev.com/download/json-%E0%B8%84%E0%B8%B7%E0%B8%AD%E0%B8%AD%E0%B8%B0%E0%B9%84%E0%B8%A3/

% \end{thebibliography}

\ifproject
\appendix
\include{chapters/appendix}

%% Display glossary (optional) -- need glossary option.
\ifglossary\glossarypage\fi

%% Display index (optional) -- need idx option.
\ifindex\indexpage\fi

\begin{biosketch}
\begin{center}
  \includegraphics[width=1.5in]{mugshot.jpg}
\end{center}
Your biosketch goes here. Make sure it sits inside
the \texttt{biosketch} environment.
\end{biosketch}
\fi % \ifproject
\end{document}
